\chapter{ws closed and ws open Maps in Topological Spaces}
\graphicspath{{Chapter4/Chapter4Figs/EPS/}{Chapter4/Chapter4Figs/}}

\section{Introduction}\label{sec4.1}

In 1982, the concept of generalized closed maps is introduced and studied by S. R, Malghan \cite{Malghan}, wg-closed maps and rwg-closed maps were introduced and studied by Nagavani \cite{Nagaveni}. Later regular closed maps, rw-closed maps and $\alpha$rw-closed maps have been introduced and studied by Long \cite{Long}, Benchalli et.all \cite{Benchalli} and R S Wali et.all \cite{Wali4} respectively. 

In section 2 of this chapter, a new class of maps called ws-closed maps is introduced and their relations with various generalized closed maps are studied. We prove that the composition of two ws-closed maps need not be ws-closed map. We also introduce ws*-closed maps, ws-open maps and ws*-open maps in topological spaces and obtain certain characterizations of these maps.

\section{ws-closed Maps and ws-open Maps}\label{sec4.2}

\begin{dfn}\label{defi4.2.1} 
A map $\TSh: \TSP \to \TSQ$ is said to be weakly semi-closed (briefly ws-closed) map if the image of every closed set in $(\TSP, \tau)$ is ws-closed in $(\TSQ, \sigma)$.
\end{dfn}

\begin{exm}\label{exam4.2.2}
Let $\TSP = \TSQ = \{1, 2, 3\}$, $\tau = \{\TSP, \phi, \{1\}, \{2\}, \{1, 2\}\}$ be a topology on $\TSP$.
$\sigma = \{\TSQ, \phi, \{1\}, \{2, 3\}\}$ be a topology on $\TSQ$ and $\ws\TSC(\TSQ) =\{\TSQ, \phi, \{2\},\{3\},\{1,3\},\break \{2, 3\}\}$. Let $\TSh: \TSP \to \TSQ$ defined by identity map, then $\TSh$ is ws-closed map.
\end{exm}

\begin{thm}\label{thm4.2.3}
Each closed map is ws-closed map, but inverse is untrue.
\end{thm}

\begin{proof}
Take up $\TSh: (\TSP,\tau)\to (\TSQ, \sigma)$ be closed map and $\TSV$ be any $\TSC(\TSP)$. Then $\TSh(\TSV)$ is closed in $\TSQ$, Seeing that each closed set is $\ws$-$\TSC(\TSP)$. Hence $\TSh(\TSV)$ is ws-closed in $\TSQ$. Therefore $\TSh$ is ws-closed.
\end{proof}

We use example \ref{exam4.2.4} to prove the inverse of theorem is untrue.

\begin{exm}\label{exam4.2.4}
Take up $\TSP = \TSQ = \{1, 2, 3\}$ $\tau = \{\TSP, \phi, \{2\},\{3\},\{1,3\},\{2, 3\}\}$ be a topology on $\TSP$. $\sigma =\{\TSQ, \phi, \{1\},\{2\},\{1, 2\}\}$ be a topology on $\TSQ$ and $\ws\TSC(\TSQ) =\{\TSQ, \phi,\{1\}, \{2\},\{3\},\{1,3\},\{2,3\}\}$. Let $\TSh: \TSP \to \TSQ$ defined by identity map, hence $\TSh$ is ws-closed map but not closed as image of closed set $\{2\}$ in $\TSP$ is $\{2\}$, is not closed in $\TSQ$.
\end{exm}

\begin{thm}\label{thm4.2.5}
If $\TSh: \TSP\to \TSQ$ is a closed map then the following holds.
\begin{enumerate}[(1)]
\item Each semi-closed is ws-closed but inverse is untrue.
\item Each $\alpha$-closed is ws-closed but inverse is untrue.
\item Each $\TSg^{\#}$-closed is ws-closed but inverse is untrue.
\item Each *g$\alpha$-closed is ws-closed but inverse is untrue.
\item Each g$\xi$*-closed is ws-closed but inverse is untrue.
\item Each $\alpha$gp-closed is ws-closed but inverse is untrue.
\item Each $\ddot{g}$-closed is ws-closed but inverse is untrue.
\item Each regular-closed is ws-closed but inverse is untrue
\item Each rb-closed is ws-closed but inverse is untrue.
\end{enumerate}
\end{thm}

\begin{proof}
Proof is obtained from the evidence that each semi (respectively, $\alpha$, g\#-, *g$\alpha$, g$\xi$*, $\ddot{g}$, regular, rb) closed set is ws-closed.
\end{proof}

We use example \ref{exam4.2.6} to prove the inverse of theorem is untrue.

\begin{exm}\label{exam4.2.6}
Take up $\TSP = \TSQ = \{1, 2, 3\}$ $\tau = \{\TSP, \phi, \{2\}, \{3\},\{1,3\},\{2, 3\}\}$ be a topology on $\TSP$. $\sigma =\{\TSQ, \phi, \{1\}, \{2\}, \{1, 2\}\}$ be a topology on $\TSQ$ and $\ws\TSC(\TSQ)=\{\TSQ, \phi,\{1\},\{2\},\{3\},\{1,3\},\{2,3\}\}$. Let $\TSh: \TSP \to \TSQ$ defined by identity map, hence $\TSh$ is ws-closed map but not closed as image of closed set $\{2\}$ in $\TSP$ is $\{2\}$, is not semi-(respectively $\alpha$, g\#, *g$\alpha$, g$\xi$*, $\ddot{g}$, regular, rb) in $\TSQ$.
\end{exm}

\begin{thm}\label{thm4.2.7}
Each g\#s-closed map is ws-closed map, but inverse is untrue.
\end{thm}

\begin{proof}
Proof is obvious, the fact that each g\#s -$\TSC(\TSP)$ is ws-$\TSC(\TSP)$.
\end{proof}

We use example \ref{exam4.2.8} to prove the inverse of theorem is untrue.

\begin{exm}\label{exam4.2.8}
Take up $\TSP=\TSQ= \{1, 2, 3, 4\}$. Let $\tau= \{\varphi, \TSP, \{1,2\}, \{3,4\}\}$ be a topology on $\TSP$ and $\sigma=\{\phi, \TSQ, \{1\}, \{2\}, \{1, 2\},\{1,2,3\}\}$ be a topology on $\TSQ$ and $\ws\TSC(\TSQ)=\{\TSQ, \phi, \{1\},\{2\},\{3\},\{4\},\{1,3\}, \{1,4\},\{2,3\},\{2,4\},\{3,4\},\{1,2,4\}, \{1,3,4\},\break \{2,3,4\}\}$, g\#s$\TSC(\TSQ)=\{\TSQ, \phi, \{2\},\{3\},\{4\},\{2,3\},\{2,4\},\{3,4\},\{2,3,4\}\}$. Let $\TSh: (\TSP, \tau )\break \to (\TSQ, \sigma)$ be a function defined by $\TSh(1)= 1$, $\TSh(2)=4$, $\TSh(3)=3$, $\TSh(4)=4$, is ws-closed map but not closed as image of closed set $\{1,2\}$ in $\TSP$ is $\{1,4\}$ is not g\#s-closed set in $\TSQ$.
\end{exm}

\begin{thm}\label{thm4.2.9}
If $\TSh: \TSP\to \TSQ$ is a map then the following holds.
\begin{enumerate}[(i)]
\item Each ws-closed map is gspr-closed map but inverse is untrue.
\item Each ws-closed map is gsp-closed map but inverse is untrue.
\item Each ws-closed map is rgb-closed map but inverse is untrue.
\end{enumerate}
\end{thm}

\begin{proof}
Proof is obtained from the evidence that each ws-$\TSC(\TSP)$ is gspr $\TSC(\TSP)$ (respectively, rgb-$\TSC(\TSP)$, gsp-$\TSC(\TSP)$).
\end{proof}

We use example \ref{exam4.2.10} to prove the inverse of theorem is untrue.

\begin{exm}\label{exam4.2.10} 
Take up $\TSP=\TSQ= \{1, 2, 3, 4\}$. Let $\tau= \{\phi, \TSP, \{1\}, \{2\}, \{1, 2\},\break \{1,2,3\}\}$ be a topology on $\TSP$ and $\sigma=\{\phi, \TSQ, \{1,2\}, \{3,4 \} \}$ be a topology on $\TSQ$ and $\ws\TSC(\TSQ)= \{\TSQ, \phi, \{1,2\},\{3,4 \}\}$, gspr$\TSC(\TSQ)=\gsp\TSC(\TSQ)=\rgb\TSC(\TSQ)=\{\TSP(\TSQ)\}$. Let $\TSh: (\TSP, \tau) \to (\TSQ,\sigma)$ defined by identity map, then $\TSh$ is gspr (respectively, rgb, gsp) closed map but not ws-closed as image of closed set $\{4\}$ in $\TSP$ is $\{4\}$ is not ws-closed in $\TSQ$.
\end{exm}

\begin{rem}\label{rem4.2.11}
The example given below shows that ws-closed maps are independent of pre, w, gp, g, $\hat{g}$, $\tilde{g}$, g$\alpha$, gw$\alpha$, pgpr, $\alpha$gp, swg-, w$\alpha\hat{g}$, $\alpha$g , $\eta$-, g*pre , mildly g, g*p, wg, rg, gpr, wgr$\alpha$, pgr$\alpha$, $\alpha$**g, rwg , $\beta$wg** , g$\alpha$**, rg$\alpha$, $\alpha$gr, gpw , $\alpha$gw, g$\alpha$*, R* closed maps,
\end{rem}

\begin{exm}\label{exam4.2.12}
Take up $\TSP=\TSQ= \{1, 2, 3, 4\}$. Let $\tau= \{\phi, \TSP, \{1\},\{1, 2\},\{1,2,3\}\}$ be a topology on $\TSP$ and $\sigma=\{\phi, \TSQ, \{1,2\}, \{3,4 \}\}$ be a topology on $\TSQ$ and $\ws\TSC(\TSQ)= \{\TSQ, \phi, \{1,2\}, \{3,4 \}\}$. Let $\TSh: \TSP \to \TSQ$ defined by, $\TSh(1) = 1$, $\TSh(2) = 2$, $\TSh(3) = 3$, $\TSh(4)=4$, then $\TSh$ is of pre, w, gp, g, $\hat{g}$, $\tilde{g}$, g$\alpha$, gw$\alpha$, pgpr, $\alpha$gp, swg-, w$\alpha\hat{g}$, $\alpha$g, $\eta$-, g*pre, mildly g, g*p, wg, rg, gpr, wgr$\alpha$, pgr$\alpha$, $\alpha$**g, rwg, $\beta$wg**, g$\alpha$**, rg$\alpha$, $\alpha$gr, gpw, $\alpha$gw, g$\alpha$*, R* closed maps, But $\TSh$ is not ws-closed map, as closed set $\{4\}$ in $\TSP$ is $\{4\}$ is not ws-closed in $\TSQ$.
\end{exm}

\begin{exm}\label{exam4.2.13}
Take up $\TSP=\TSQ= \{1, 2, 3, 4\}$. Let $\tau= \{\phi, \TSP, \{1,2\}, \{3,,4 \}\}$ be a topology on $\TSP$ and $\sigma=\{\phi, \TSQ, \{1\}, \{2\}, \{1,2\},\{1,2,3\}\}$ be a topology on $\TSQ$ and $\ws\TSC(\TSQ)=\{\phi, \TSQ, \{1\},\{2\},\{3\},\{4\},\{1,3\}, \{1,4\},\{2,3\},\{2,4\},\{3,4\},\{1,2,4\},\break \{1,3,4\}, \{2,3,4\}\}$. Take up $\TSh: \TSP \to \TSQ$ defined by, $\TSh(1) = 1$, $\TSh(2) = 1$, $\TSh(3) = 3$, $\TSh(4)=4$, then $\TSh$ is ws-closed, but not a pre(respectively, w, gp, g, $\hat{g}$, $\tilde{g}$, g$\alpha$, gw$\alpha$, pgpr, $\alpha$gp, swg, w$\alpha\hat{g}$, $\alpha$g, $\eta$, g*pre, mildly g, g*p, wg, rg, gpr, wgr$\alpha$, pgr$\alpha$, $\alpha$**g, rwg, $\beta$wg**, g$\alpha$**, rg$\alpha$, $\alpha$gr, gpw, $\alpha$gw, g$\alpha$*, R*) closed maps as $\TSC(\TSP)$ $\{\TSa , \TSb\}$ in $\TSP$ is $\{1\}$ is not pre(respectively, w, gp, g, $\hat{g}$, $\tilde{g}$, g$\alpha$, gw$\alpha$, pgpr, $\alpha$gp, swg, w$\alpha\hat{g}$, $\alpha$g, $\eta$, g*pre, mildly g, g*p, wg, rg, gpr, wgr$\alpha$, pgr$\alpha$, $\alpha$**g, rwg, $\beta$wg**, g$\alpha$**, rg$\alpha$, $\alpha$gr, gpw, $\alpha$gw, g$\alpha$*, R* ) closed set in $\TSQ$.
\end{exm}

\begin{rem}\label{rem4.2.14}
From the above discussions and known facts, the relation between ws-closed map and some existing closed maps in topological space is shown in the following figure
\begin{figure}[!ht]
\centering
\includegraphics[scale=.75]{1.eps}
\end{figure}
\end{rem}

\begin{thm}\label{thm4.2.15}
Map $\TSh: \TSP \to \TSQ$ is contra regular closed and gspr-closed map then $\TSh$ is ws-closed map in $\TSP$.
\end{thm}

\begin{proof}
Take up $\TSM$ as a closed set $\in (\TSP, \tau)$. Then $\TSh (\TSM)$ will be open and gspr-closed in $\TSQ$. By Results 1.3.3 $\TSh (\TSM)$ is ws-closed in $(\TSQ, \sigma)$. Therefore $\TSh$ is ws-closed map.
\end{proof}

\begin{thm}\label{thm4.2.16}
A map $\TSh: \TSP \to \TSQ$ is contra closed and rgb-closed map then h is ws-closed map in $\TSP$.
\end{thm}

\begin{proof}
Take up $\TSM$ as a closed set $\in (\TSP, \tau)$. Then $\TSh (\TSM)$ will be open and rgb - closed in $\TSQ$. By Results 1.3.3 $\TSh (\TSM)$ is ws-closed in $(\TSQ, \sigma)$. Therefore $\TSh$ is ws-closed map.
\end{proof}

\begin{thm}\label{thm4.2.17}
Map $\TSh: \TSP \to \TSQ$ is contra closed and $\alpha$g-closed map then h is ws-closed map.
\end{thm}

\begin{proof}
Take up $\TSM$ as a closed set $\in (\TSP, \tau)$. Then $\TSh (\TSM)$ will be open and $\alpha$g-closed in $\TSQ$. By Results 1.3.3 $\TSh (\TSM)$ is ws-closed in $(\TSQ, \sigma)$. Therefore $\TSh$ is ws-closed map.
\end{proof}

\begin{thm}\label{thm4.2.18}
Map $\TSh: \TSP \to \TSQ$ is contra closed and rgb-closed map then $\TSh$ is ws-closed map in $\TSP$.
\end{thm}

\begin{proof}
Take up $\TSM$ as a closed set $\in (\TSP, \tau)$. Then $\TSh (\TSM)$ will be open and rgb - closed in $\TSQ$. By results 1.3.3 $\TSh(\TSM)$ is ws-closed in $(\TSQ, \sigma)$. Therefore $\TSh$ is ws-closed map.
\end{proof}

\begin{thm}\label{thm4.2.19}
Map $\TSh: \TSP \to \TSQ$ is contra semi closed and swg*-closed map then $\TSh$ is ws-closed map in $\TSP$.
\end{thm}

\begin{proof}
Take up $\TSM$ as a closed set $\in (\TSP, \tau)$. Then $\TSh (\TSM)$ will be semi open and swg*-closed in $\TSQ$. By Results 1.3.3 $\TSh(\TSM)$ is ws-closed in $(\TSQ, \sigma)$. Therefore $\TSh$ is ws-closed map.
\end{proof}

\begin{thm}\label{thm4.2.20}
Map $\TSh: \TSP \to \TSQ$ is contra semi closed and swg-closed map then $\TSh$ is ws-closed map in $\TSP$.
\end{thm}

\begin{proof}
Take up $\TSM$ as a closed set $\in (\TSP, \tau)$. Then $\TSh (\TSM)$ will be semi open and swg-closed in $\TSQ$. By Results 1.3.3 $\TSh (\TSM)$ is ws-closed in $(\TSQ, \sigma)$. Therefore $\TSh$ is ws-closed map.
\end{proof}

\begin{thm}\label{thm4.2.21}
Map $\TSh: \TSP \to \TSQ$ is contra is semi closed and sg-closed map then $\TSh$ is ws-closed map in $\TSP$.
\end{thm}

\begin{proof}
Take up $\TSM$ as a closed set $\in(\TSP, \tau)$. Then $\TSh (\TSM)$ will be semi open and sg-closed in $\TSQ$. By Results 1.3.3 $\TSh (\TSM)$ is ws-closed in $(\TSQ, \sigma)$. Therefore $\TSh$ is ws-closed map.
\end{proof}

\begin{thm}\label{thm4.2.22}
Map $\TSh: \TSP \to \TSQ$ is contra is semi closed and sgb-closed map then $\TSh$ is ws-closed map in $\TSP$.
\end{thm}

\begin{proof}
Take up $\TSM$ as a closed set $\in (\TSP, \tau)$. Then $\TSh(\TSM)$ will be semi open and sgb - closed in $\TSQ$. By Results 1.3.3 $\TSh (\TSM)$ is ws-closed in $(\TSQ, \sigma)$. Therefore $\TSh$ is ws-closed map.
\end{proof}

\begin{thm}\label{thm4.2.23}
Map $\TSh: \TSP \to \TSQ$ is contra is semi closed and $\alpha$gs-closed map then $\TSh$ is ws-closed map in $\TSP$.
\end{thm}

\begin{proof}
Take up $\TSM$ as a closed set $\in (\TSP, \tau)$. Then $\TSh (\TSM)$ will be semi open and $\alpha$gs-closed in $\TSQ$. By Results 1.3.3 $\TSh (\TSM)$ is ws-closed in $(\TSQ, \sigma)$. Therefore $\TSh$ is ws-closed map.
\end{proof}

\begin{thm}\label{thm4.2.24}
Map $\TSh: \TSP \to \TSQ$ is contra is $\beta$-closed and $\beta$wg*-closed map then $\TSh$ is ws-closed map in $\TSP$.
\end{thm}

\begin{proof}
Take up $\TSM$ as a closed set $\in (\TSP, \tau)$. Then $\TSh (\TSM)$ is $\beta$-open and $\beta$wg*-closed in $\TSQ$. By Results 1.3.3 $\TSh (\TSM)$ is ws-closed in $(\TSQ, \sigma)$. Therefore $\TSh$ is ws-closed map.
\end{proof}

\begin{thm}\label{thm4.2.25}
Map $\TSh: \TSP \to \TSQ$ is contra is both closed and $\TSg$-closed map then $\TSh$ is ws-closed map in $\TSP$.
\end{thm}

\begin{proof}
Take up $\TSM$ as a closed set $\in (\TSP, \tau)$. Then $\TSh (\TSM)$ will be open and $\TSg$-closed in $\TSQ$. By Results 1.3.3 $\TSh (\TSM)$ is ws-closed in $(\TSQ, \sigma)$. Therefore $\TSh$ is ws-closed map.
\end{proof}

\begin{thm}\label{thm4.2.26}
Map $\TSh: \TSP \to \TSQ$ is contra is regular semi closed and rw-closed map then $\TSh$ is ws-closed map in $\TSP$.
\end{thm}

\begin{proof}
Take up $\TSM$ as a closed set $\in(\TSP, \tau)$. Then $\TSh (\TSM)$ will be semi open and rw-closed in $\TSQ$. By Results 1.3.3 $\TSh (\TSM)$ is ws-closed in $(\TSQ, \sigma)$. Therefore $\TSh$ is ws-closed map.
\end{proof}

\begin{thm}\label{thm4.2.27}
Map $\TSh: \TSP \to \TSQ$ is contra is regular semi closed and R*-closed map then $\TSh$ is ws-closed map in $\TSP$.
\end{thm}

\begin{proof}
Take up $\TSM$ as a closed set $\in (\TSP, \tau)$. Then $\TSh (\TSM)$ will be semi open and R*-closed in $\TSQ$. By Results 1.3.3 $\TSh (\TSM)$ is ws-closed in $(\TSQ, \sigma)$. Therefore $\TSh$ is ws-closed map.
\end{proof}

\begin{thm}\label{thm4.2.28}
A mapping $\TSh: \TSP \to \TSQ$ be ws-closed, then $\wscl (\TSh(\clrD))\subseteq \TSh(\cl (\clrD)) \forall$ subset $\clrD$ of $(\TSP, \tau)$.
\end{thm}

\begin{proof}
Pretend that $\TSh$ is ws-closed and $\clrD \subseteq \TSP$. Then $\cl(\clrD)$ is closed in $\TSP$ and so $\TSh(\cl(\clrD))$ is ws-closed in $(\TSQ, \sigma)$. We have $\TSh(\clrD) \subseteq \TSh(\cl(\clrD))$, $\wscl(\TSh(\clrD)) \subseteq \wscl(\TSh(\cl(\clrD)))\break \to$ (1). Since $\TSh(\cl(\clrD))$ is ws-closed in $(\TSQ, \sigma)$, $\wscl(\TSh(\cl(\clrD))) = \TSh (\cl (\clrD)) \to$ (2), From (1) and (2), we have wscl $(\TSh(\clrD)) \subseteq \TSh(\cl(\clrD))$ for each subset $\clrD$ of $(\TSP, \tau)$.
\end{proof}

\begin{coro}\label{coro4.2.29}
A mapping $\TSh: \TSP \to \TSQ$ be ws-closed, then the image $\TSh(\clrD)$ of closed set $\clrD$ in $(\TSP, \tau)$ is $\tau_{\ws}$-closed in $(\TSQ, \sigma)$.
\end{coro}

\begin{proof}
Take up closed set $\clrD \in (\TSP, \tau)$. Since $\TSh$ is ws-closed, wscl $(\TSh (\clrD)) \subseteq \TSh (\cl (\clrD)) \to$ (1). Also $\cl (\clrD) = \clrD$, as $\clrD$ is a closed set and so $\TSh (\cl(\clrD)) = \TSh(\clrD) \to$ (2). From (1) and (2), we have $\wscl (\TSh (\clrD)) \subseteq \TSh (\clrD)$. We know that $\TSh (\clrD) \subseteq \ws \cl (\TSh (\clrD))$ and so $\wscl (\TSh (\clrD)) =\TSh (\clrD)$. Therefore $\TSh (\clrD)$ is $\tau_{\ws}$-closed in $(\TSQ, \sigma)$.
\end{proof}

\begin{thm}\label{thm4.2.30}
Take up $(\TSP, \tau)$ be any topological spaces and $\TSQ$ be a topological space where $\wscl (\clrD) = \alpha - \cl (\clrD) \forall \clrD$ of $\TSQ$, and $\TSh: \TSP \to \TSQ$ be a map, and then given statements are one and the same..
\begin{enumerate}[(i)]
\item $\TSh$ is ws-closed map.
\item $\wscl (\TSh (\clrD)) \subseteq \TSh (\cl (\clrD)) \forall$ subset $\clrD$ of $(\TSP, \tau)$.
\end{enumerate}
\end{thm}

\begin{proof}
(i) $\Rightarrow$ (ii) obiuosly from Theorem \ref{thm4.2.28}.

(ii) $\Rightarrow$ (i) Take up closed set $\clrD \in (\TSP, \tau)$. Then $\clrD=\cl (\clrD)$ and so $\TSh(\clrD) = h(\cl (\clrD)) \supseteq \wscl (\TSh (\clrD))$ by hypothesis. We have $\TSh (\clrD) \subseteq \wscl (\TSh (\clrD))$, Therefore $\TSh (\clrD) = \wscl (\TSh (\clrD))$. Also $\TSh (\clrD) = \wscl (\TSh (\clrD)) = \alpha-\cl (\TSh (\clrD))$, by hypothesis. That is $\TSh (\clrD) = \alpha-\cl (\TSh (\clrD))$ and so $\TSh (\clrD)$ is $\alpha$-closed in $(\TSQ, \sigma)$. Thus $\TSh (\clrD)$ is ws-closed in $\TSQ$ and therefore $\TSh$ is ws-closed map.
\end{proof}

\begin{thm}\label{thm4.2.31}
Map $\TSh : \TSP \to \TSQ$ be ws-closed iff $\forall$ subset $\TSS$ of $\TSQ$ and each open set $\TSM$ containing $h^{-1}(\TSS) \subseteq \TSM$, there is a ws-open set $\TSV$ of $\TSQ \ni: S \subseteq \TSV$ and $h^{-1} (\TSV) \subseteq \TSM$.
\end{thm}

\begin{proof}
Pretend $\TSh$ is ws-closed. Assume $\TSS \subseteq \TSQ$ and $\TSM$ is an open set $\in \TSP$, $\ni : h^{-1} (\TSS) \subseteq \TSM$.  Now $\TSP$-$\TSM$ will be closed in $\TSP$. Seeing that $\TSh$ is ws-closed, $\TSh(\TSP\text{-}\TSM)$ is ws closed $\in\TSQ$. Then $\TSV = \TSQ - \TSh (\TSP\text{-}\TSM)$ will be ws-open set $\in\TSQ$. and it is known that $h^{-1} (\TSS) \subseteq \TSM$ implies $\TSS \subseteq \TSV$ and $h^{-1}(\TSV) =\TSP - \TSh^{-1} (\TSh(\TSP\text{-}\TSM)) \subseteq \TSP - (\TSP\text{-}\TSM) = \TSM$. That is $h^{-1}(\TSV)\subseteq  \TSM$.

Inversely,

Pretend $\TSH$ is $\TSC(\TSP)$. Then $h^{-1}((\TSh (\TSH))^{\text{c}}) \subseteq \TSH^{\text{c}}$ and $\TSH^{\text{c}}$ will be open set $\in\TSP$. From hypothesis, $\exists$ an ws-open set $\TSV \in \TSQ \ni: \TSh(\TSH)^{\text{c}} \subseteq \TSV$ and $h^{-1}(\TSV) \subseteq \TSH^{\text{c}}$ and so $\TSH \subseteq (h^{-1} (\TSV))^{\text{c}}$. Hence $\TSV^{\text{c}} \subseteq \TSh(\TSH) \subseteq \TSh(((h^{-1} (\TSV))^{\text{c}} \subseteq \TSV^{\text{c}}$ which implies $\TSh(\TSH) = \TSV^{\text{c}}$. Seeing that $\TSV^{\text{c}}$ is ws-closed, $\TSh (\TSH)$ is ws-closed. Thus $\TSh (\TSH)$ is ws-closed $\in \TSQ$ and therefore $\TSh$ is ws-closed map.
\end{proof}

\begin{rem}\label{rem4.2.32}
As usual, The composition of two ws-closed maps need not be ws-closed map and it is proved from the example given below.
\end{rem}

\begin{exm}\label{exam4.2.33}
Take up $\TSP = \TSQ = \TSR = \{1,2,3\}$, $\tau=\{\TSP, \phi, \{1\}, \{2, 3\}\}$ be a topology on $\TSP$, $\sigma = \{\TSQ, \phi, \{1\},\{2\},\{1, 2\}\}$ be a topology on $\TSQ$ and $\eta = \{\TSR, \phi, \{1\}, \{2\},\break \{1, 2\}, \{1, 3\}\}$ be a topology on $\TSR$. Define $\TSh: \TSP \to \TSQ$ defined by $\TSh(1)=1$, $\TSh(2)=2$, $\TSh(3)=3$ and $\TSg: \TSQ\to (\TSR)$ are the identity maps. Then $\TSh$ and $\TSg$ are ws-closed maps, but their composition $\TSg\circ \TSh : (\TSP) \to (\TSR)$ will not be ws-closed map, because $\TSH = \{1\}$ is closed in $(\TSP)$, but $\TSg \circ \TSh (\TSH) =\TSg \circ \TSh (\{1\}) = \TSg(\TSh(\{1\})) = \TSg(\{1\}) = \{1\}$ is not ws-closed in $(\TSR)$.
\end{exm}

\begin{thm}\label{thm4.2.34}
If $\TSh: \TSP \to \TSQ$ is closed map and $\TSg: \TSQ\to (\TSR)$ is ws-closed map, then the composition $\TSg \circ \TSh: (\TSP, \tau) \to (\TSR, \eta)$ is ws-closed map.
\end{thm}

\begin{proof}
Take up closed set $\TSF \in(\TSP)$. Seeing that $\TSh$ is closed map, $\TSh (\TSF)$ will be closed in $(\TSQ)$. Also $\TSg$ is ws-closed map, $\TSg (\TSh(\TSF))$ will be ws-closed in $(\TSR)$. That is $\TSg \circ \TSh (\TSF) = \TSg (\TSh (\TSF))$ is ws-closed and henceforth $\TSg \circ \TSh$ is ws-closed map.
\end{proof}

\begin{thm}\label{thm4.2.35}
If $\TSh: \TSP \to \TSQ$ and $\TSg: \TSQ\to (\TSR)$ is ws-closed maps and $\TSQ$ be a $\TST_{\text{ws}}$-space then $\TSg \circ \TSh: (\TSP, \tau) \to (\TSR, \eta)$ is ws-closed map.
\end{thm}

\begin{proof}
Take up closed set $\clrD \in (\TSP, \tau)$. Seeing that $\TSh$ is ws-closed, $\TSh (\clrD)$ is ws-closed in $\TSQ$. From by hypothesis, $\TSh (\clrD)$ is closed. Also $\TSg$ is ws-closed, $\TSg (\TSh (\clrD))$ will be ws-closed in $\TSR$ and $\TSg(\TSh(\clrD)) = \TSg \circ \TSh(\clrD)$. Henceforth $\TSg \circ \TSh$ is ws-closed map.
\end{proof}

\begin{thm}\label{thm4.2.36}
If $\TSh: \TSP \to \TSQ$ is $\TSg$-closed, $\TSg: \TSQ \to \TSR$, be ws-closed and $\TSQ$ is T1/2-space then their composition $\TSg \circ \TSh: \TSP \to \TSR$ is ws-closed map.
\end{thm}

\begin{proof}
Take up closed set $\clrD \in (\TSP, \tau)$. Seeing that $\TSh$ is $\TSg$-closed, $\TSh (\clrD)$ will be $\TSg$-closed in $(\TSQ)$. Seeing that $\TSQ$ is T1/2-space, $\TSh(\clrD)$ is closed $\in (\TSQ)$. Since $\TSg$ is ws-closed, $\TSg (\TSh(\clrD))$ is ws-closed $\in \TSR$ and $\TSg(\TSh(\clrD)) = \TSg \circ \TSh(\clrD)$. Henceforth $\TSg \circ \TSh$ is ws-closed map.
\end{proof}

\section{ws-Open Maps}

\begin{dfn}\label{defi4.2.37}
A map $\TSh: \TSP \to \TSQ$ is called a ws-open map if the image $\TSh (\clrD)$ is ws-open in $\TSQ \ \forall$ open set $\clrD$ in $(\TSP)$. 

From the definitions we obtain these following results.
\end{dfn}

\begin{thm}\label{thm4.2.38}
\begin{enumerate}[(i)]
\item Each open map is ws-open but inverse is untrue.
\item Each semi-open map is ws-open but inverse is untrue.
\item Each $\alpha$-open map is ws-open but inverse is untrue.
\item Each g\#-open map is ws-open but inverse is untrue.
\item Each *g$\alpha$-open map is ws-open but inverse is untrue.
\item Each g$\xi$*-open map is ws-open but inverse is untrue.
\item Each $\alpha$gp-open map is ws-open but inverse is untrue.
\item Each $\ddot{g}$-open map is ws-open but inverse is untrue.
\item Each rb-regular open map is ws-open but inverse is untrue
\item Each g\#s open map is ws-open but inverse is untrue.
\item Each ws-open map is gspr-open but inverse is untrue.
\item Each ws-open map is gsp-open but inverse is untrue.
\item Each ws-open map is rgb-open but inverse is untrue.
\end{enumerate}
\end{thm}

\begin{thm}\label{thm4.2.39}
For each bijection map $\TSh: \TSP \to \TSQ$ following results are one and the same.
\begin{enumerate}[(1)]
\item $h^{-1} : \TSQ \to \TSP$ is ws-continuous.
\item $\TSh$ is ws-open map
\item $\TSh$ is ws-closed map.
\end{enumerate}
\end{thm}

\begin{proof}
(1) $\Rightarrow$ (2) Take up a open set $\TSM \in \TSP$. By hypothesis $(h^{-1})^{-1}(\TSM) = \TSh (\TSM)$ is ws-open $\in \TSQ$ and so $\TSh$ is ws-open.

(2) $\Rightarrow$ (3) Take up closed set $\TSH\in \TSP$. Then $\TSH^{\text{c}}$ is open set $\in \TSP$. By hypothesis, $\TSh (\TSH^{\text{c}})$ is ws-open $\in \TSQ$. Also $\TSh (\TSH^{\text{c}}) = \TSh (\TSH)^{\text{c}}$ is ws-open $\in \TSQ$ and therefore $\TSh(\TSH)$ is ws-closed in $\TSQ$. Hence $\TSh$ is ws-closed.

(3) $\Rightarrow$ (4) Take up closed set $\TSH \in \TSP$. By hypothesis, $\TSh(\TSH)$ is ws-closed $\in \TSQ$. But $\TSh (\TSH) = (h^{-1})^{-1}(\TSH)$ and henceforth $h^{-1}$ is continuous.
\end{proof}

\begin{thm}\label{thm4.2.40}
A Map $\TSh: \TSP\to\TSQ$ is ws-open, then $\TSh (\text{int}(\clrD)) \subseteq$ ws-int $(\TSh(\clrD))$ for $\forall$ subset $\clrD\in \TSP$.
\end{thm}

\begin{proof}
Take up $\TSh: \TSP\to \TSQ$ be an open map and $\clrD$ is any subset $\in\TSP$. Hence $\text{int}(\clrD)$ is open $\in \TSP$ and so $\TSh(\text{int}(\clrD))$ is ws-open $\in\TSQ$. We have $\TSh(\text{int}(\clrD)) \subseteq \TSh(\clrD)$. Henceforth $\TSh(\text{int} (\clrD)) \subseteq$ ws-int $(\TSh(\clrD))$.
\end{proof}

\begin{thm}\label{thm4.2.41}
A map $\TSh: \TSP\to \TSQ$ is ws-open, then for each neighbourhood $\TSM$ of $\TSx \in \TSP \exists$ a ws-neighbourhood $\TSW$ and $\TSh(\TSx) \in \TSQ$, $\ni:\TSW \subseteq \TSh(\TSM)$.
\end{thm}

\begin{proof}
Take up $\TSh: \TSP\to \TSQ$ is an ws-open map. Assume $\TSx\in \TSP$ and $\TSM$ be an arbitrary neighbourhood of $\TSx\in \TSP$. Then $\exists$ an open set $\TSG \in \TSP$, $\ni: \TSx \in \TSG \subseteq \TSM$. Also $\TSh(\TSx)\in \TSh (\TSG) \subseteq \TSh (\TSM)$ and $\TSh (\TSG)$ is ws-open set $\in \TSQ$, as $\TSh$ is an ws-open map. And also by the result $\TSh (\TSG)$ is ws-neighbourhood of each of its points. Taking $\TSh (\TSG) = \TSW$, $\TSW$ is an ws-neighbourhood of $\TSh(\TSx)$ in $\TSQ \ni: \TSW \subseteq \TSh (\TSM)$.
\end{proof}

\begin{dfn}\label{defi4.2.42}
A map $\TSh: \TSP\to \TSQ$ is termed as ws*-closed map if the image $\TSh(\clrD)$ is ws-closed $\in \TSQ$ for each ws-closed set $\clrD \in(\TSP, \tau)$.
\end{dfn}

\begin{thm}\label{thm4.2.43}
Each ws*-closed map is ws-closed map but inverse is untrue.
\end{thm}

\begin{proof}
Proof is obtained from the previous results, definations and evidence that each closed set is ws-closed.
\end{proof}

\begin{exm}\label{exam4.2.44}
Take up $\TSP = \TSQ = \{1, 2, 3\}$, $\tau = \{\TSP, \phi, \{1\}, \{2\}, \{1, 2\}\}$ be a topology on $\TSP$ and $\sigma = \{\TSQ, \phi, \{1\}, \{2\},\{1, 2\},\{1,3\}\}$ be a topology on $\TSQ$ and $\TSh: \TSP \to \TSQ$ be the identity map. Thus $\TSh$ is ws-closed map but not ws*-closed map. Since $\{1\}$ is ws-closed in $\TSP$ but its image under $\TSh$ is $\{1\}$, is not ws-closed in $\TSQ$.
\end{exm}

\begin{thm}\label{thm4.2.45}
If $\TSh: \TSP \to \TSQ$ and $\TSg: \TSQ\to (\TSR)$ are ws*-closed maps, then their composition $\TSg \circ \TSh: (\TSP) \to (\TSR)$ is also ws*-closed.
\end{thm}

\begin{proof}
Take up a ws-closed set $\TSH \in \TSP$. Seeing that $\TSh$ is ws*-closed map, $\TSh (\TSH)$ is ws-closed $\in \TSQ$. Seeing that ws*- closed map, $\TSg (\TSh (\TSH))$ is ws-closed $\in (\TSR)$. Henceforth $\TSg \circ \TSh$ is ws*-closed map.
\end{proof}

\begin{dfn}\label{defi4.2.46}
A map $\TSh: \TSP \to \TSQ$ is termed as ws*-open map if the image $\TSh (\clrD)$ is ws-open set in $\TSQ$ for each ws-open set $\clrD$ in $\TSP$.
\end{dfn}

\begin{rem}\label{rem4.2.47}
Seeing that each open set is a ws-open set, we have each ws*-open map is ws-open map. As usual the inverse is untrue as seen from example given below.
\end{rem}

\begin{exm}\label{exam4.2.48}
Take up $\TSP = \TSQ = \{1, 2, 3\}$, $\tau = \{\TSP, \phi, \{1\}, \{2\}, \{1, 2\}\}$ and $\sigma =\{\TSQ, \phi, \{1\}, \{2\}, \{1, 2\},\{1,3\}\}$ and $\TSh: \TSP \to \TSQ$ be the identity map. Then $\TSh$ is ws-open map but not ws*-open map. Seeing that $\{2, 3\}$ is ws-open set $\in\TSP$, but its image under $\TSh$ is $\{2, 3\}$ is not ws-open in $\TSQ$.
\end{exm}

\begin{thm}\label{4.2.49}
If $\TSh: \TSP \to \TSQ$ and $\TSg: \TSQ \to \TSR$, are ws*-open maps, then their composition $\TSg \circ \TSh: \TSP \to \TSR$, is also ws*-open.
\end{thm}

\begin{proof}
Proof is same as Theorem \ref{thm4.2.45}.
\end{proof}

\begin{thm}\label{thm4.2.50} 
For each bijection map $\TSh: \TSP \to \TSQ$ then given results are one and the same:
\begin{enumerate}[(i)]
\item $h^{-1}: \TSQ\to (\TSP, \tau)$ is ws irresolute.
\item $\TSh$ is ws*-open map
\item $\TSh$ is ws*-closed map.
\end{enumerate}
\end{thm}

\begin{proof}
Proof is similar to Theorem \ref{thm4.2.43}.
\end{proof}
